\bigskip
\begin{center}\textls[120]{b) Die lemniskatischen Funktionen in den Jahren 1796—1799.}\end{center}
\indent In der Tagebuchnotiz 33. vom September 1796 schreibt \textsc{Gauss} die Entwicklung
der Umkehrungsfunktion des Integrals \(z=\displaystyle\int\limits_0^x\frac{dx}{\sqrt{1-x^n}}\) nach Potenzen
von \(z\) in der Form
\[
z-\frac{1\,.\,z^n}{2(n+1)}A+\frac{(n-1)z^n}{4(2n+1)}B
-\frac{(n^2-n-1)z^n}{2(n+1)(3n+1)}C+\ldots,
\]
eine Schreibweise, die auf \textsc{Stirling} hinweist, der sie auf Seite 3 seines Buches
\textit{Methodus Differentialis}\footnote{\textit{Methodus differentialis sive tractatus de summatione et interpolatione serierum infinitarum} auctore \textsc{Iacobo Stirling} R. S. S. Londini MDCCXXX. Hier ist zu bemerken, dass für \textsc{Stirling} als orthodoxen Newtonianer Methodus Differentialis soviel heisst wie die Differenzenrechnung, für Differentialrechnung sagt er Methodus Fluxionum, z. B. heisst es auf S. 36 des angeführten Werks: »Der Summation der Reihen in der Differentialmethode entspricht die Quadratur der Kurven in der Fluxionsmethode«.}
auf \textsc{Newton} zurückführt und wie folgt erklärt: »Die
Anfangsglieder einer Reihe bezeichne ich durch die Anfangsbuchstaben des
Alphabets \(A,B,C,D\) etc., \(A\) ist das erste, \(B\) das zweite, \(C\) das dritte und
so fort.« Wenn nun \textsc{Gauss} in der Tagebuchnotiz 32. (Werke X, 1, S. 502)
und im Art. 5 der \textit{Exercitationes} (ebenda, S. 140) die Umkehrungsfunktion des
Falles \(n=3\), und im Art. 4 der \textit{Exercitationes} (ebenda, S. 141) die des Falles
\(n=4\) besonders hinschreibt, so wird man daraus nur als Hervorheben der nach
\(n=2\) einfachsten Fälle anzusehen haben. Erst in der Tagebuchnotiz 51.\par
\clearpage\SourceLeaf{116}\ArithmeticHead{31}{R}{III b. LEMNISKATISCHE FUNKTIONEN 1796—1799.}
\noindent vom 8. Januar 1797 (»Die von \(\displaystyle\int\frac{dx}{\sqrt{1-x^4}}\) abhängende lemniskatische Kurve habe
ich zu untersuchen angefangen«) erscheint der Fall \(n=4\) als Gegenstand, der
einer besondern Behandlung würdig ist, und wir werden durch die Tagebuchnotiz
50. vom 7. Januar 1797 auf \textsc{Euler} als die Quelle hingewiesen, der die
Anregung zu dieser besondern Behandlung entstammt. Die beiden dort auftretenden
lemniskatischen Integrale mit der trigonometrischen Transformation
finden sich nämlich u. a. im Kapitel VIII des ersten Bandes der \textit{Inst. calc.
integralis}, und dass \textsc{Gauss} zu jener Zeit mit diesem Werke beschäftigt war,
zeigen die Tagebuchnotizen 52., 53., 54., 59., die direkt an \textsc{Euler} anknüpfen
(vergl. die Bemerkungen Werke X, 1, S. 509ff.). Auch in den Leistenotizen,
die Werke X, 1, S. 145ff. abgedruckt sind, und die aus derselben Zeit stammen,
finden wir besonders in den Artikeln 1, 3, 4 Anknüpfung an \textsc{Euler} u. z.
an dessen Abhandlung \textit{De miris proprietatibus curvae elasticae}\footnote{1782—86, \textsc{Euleri} Opera ser. I, vol. 21, S. 91, vergl. auch Werke X, 1, S. 149.}.
Im Art. 2 wird
die Reihenentwicklung für das Integral \(\displaystyle\int\frac{dx}{\sqrt{1-x^4}}\) und die für seine Umkehrung
rekapituliert, aber die in den Artikeln 1, 3 und 8 enthaltenen numerischen
Werte für die kompletten Integrale
\[
A=\int\limits_0^1\frac{dx}{\sqrt{1-x^4}}\mathpunct{\vcenter{\hbox{,}}}\quad
B=\int\limits_0^1\frac{x^2\,dx}{\sqrt{1-x^4}}
\]
erfordern eine besonders aufmerksame Betrachtung. — Dass der im Art. 1
angegebene Wert \(A=1{,}311031\) von \textsc{Euler} herrührt (a. a. O., S. 104) wird
zwar nicht hier, wohl aber in einer aus dem Jahre 1798 stammenden Aufzeichnung
der Scheda Aa (Werke III, S. 413) ausdrücklich bemerkt; der ebenfalls
im Art. 1 angegebene \textsc{Stirling}sche Wert, bei dem \textsc{Gauss} ausdrücklich
auf die \textit{Methodus Differentialis} Bezug nimmt (der Wert wird daselbst S. 58
aus einer Reihenentwicklung berechnet), zeigt, dass \textsc{Eulers} Wert zu gross ist.
Da nach \textsc{Euler} (a. a. O., S. 106) \(AB=\frac14\pi\) ist, erscheint der im Art. 3 gegebene
\textsc{Euler}sche Wert von \(B\) entsprechend zu klein; \textsc{Stirling} hat (a. a. O., S. 57)
auch den Wert von \(B\) auf 17 Dezimalstellen. In der Scheda Aa (Werke III,
S. 413), also 1798, verifiziert \textsc{Gauss} den \textsc{Stirling}schen Wert von \(A\) »mit Anwendung
der Formel \(\operatorname{arc\,sin\,lemn}\frac7{23}+2\operatorname{arc\,sin\,lemn}\frac12\)«, also mit Hilfe des\par
\clearpage\SourceLeaf{117}\ArithmeticHead{32}{L}{L. SCHLESINGER, ÜBER GAUSS’ FUNKTIONENTHEORETISCHE ARBEITEN.}
\noindent Additionstheorems, und in der Abhandlung \textit{circa seriem} (1812, Werke III,
S. 150) gibt er einen noch genaueren Wert, den er »durch eine auf einen
besonderen Kunstgriff gegründete Rechnung« gefunden hat. Welche grosse
Bedeutung gerade der Wert von \(B\) oder, mit der von \textsc{Gauss} später angewandten
Bezeichnung \(\frac12\varpi=\displaystyle\int\limits_0^1\frac{dx}{\sqrt{1-x^4}}\mathpunct{\vcenter{\hbox{,}}}\)
der Wert \(2B=\frac{\pi}{\varpi}\) für die Entwicklung
des \textsc{Gauss}schen Gedankenganges gehabt hat, werden wir weiter unten ausführlich
zu zeigen haben. Hier bemerken wir nur noch, dass im Art. 8 (Werke X, 1,
S. 149) die \(A,B\), bezw. lineare Kombinationen dieser Werte, als die beiden ersten
Entwicklungskoeffizienten von \((1+\sin^2\varphi)^{-\frac12}\) nach den Kosinus der Vielfachen
von \(\varphi\) auftreten, eine Entwicklung, die sich auch im \textit{Calculus integralis} von
\textsc{Euler} findet (Opera ser. I, vol. 11, S. 98). \textsc{Gauss} benutzt dort die erhaltenen
Reihen auch zur numerischen Verifikation der Gleichung \(\frac{2A}{\pi}=\frac1{2B}\); eine
geschicktere Ausführung derselben Rechnung findet sich auf der Rückseite
des Titelblattes des \textit{Leiste} (Werke X, 1, S. 169, Art. 5, vergl. die Bemerkung
S. 171). Als wesentliches Moment erscheint in diesen Aufzeichnungen die
Anwendung des \textsc{Euler}schen Additionstheorems für das Integral erster Gattung,
das dann im Art. 4 (Werke X, 1, S. 147) zu den Additions- und Multiplikationsformeln
für die Umkehrungsfunktion ausgestaltet wird, ferner aber der
Übergang von den bloss formalen Methoden des Rechnens mit Reihen zu
numerischen Bestimmungen, die als eine Art Zwischenstufe für den Fortschritt
zu wirklich quantitativen Methoden (Konvergenzuntersuchungen, Restabschätzungen)
angesehen werden können.

\textsc{Gauss} hat sich wohl darum dazu entschlossen, das lemniskatische Integral
zum Gegenstande eingehender Untersuchungen zu machen, weil er in ihm
auf Grund seines Ursprungs aus der Rektifikation der Lemniskate die genuine
Analogie des trigonometrischen Integrals zu finden hoffte; dass das Integral,
das den Bogen der Ellipse darstellt, keine Analogie mit dem den Kreisbogen
liefernden trigonometrischen Integrale darbietet, musste ihm schon aus der
Form des Additionstheorems für die Integrale erster Gattung deutlich werden.
Aus dem Additionstheorem erschliesst er nun des weitern, dass die Umkehrfunktion
des lemniskatischen Integrals eine eindeutige und periodische Funktion
ist.
